\documentclass[11pt]{article}
\usepackage[T1]{fontenc}
\usepackage[utf8]{inputenc}
\usepackage{amsmath,amssymb,amsthm,mathtools}
\usepackage{geometry,booktabs,array,longtable}
\usepackage[hidelinks]{hyperref}
\usepackage{url}
\geometry{margin=28mm}
\setlength{\emergencystretch}{2em}
\newtheorem{theorem}{Theorem}[section]
\newtheorem{proposition}[theorem]{Proposition}
\newtheorem{lemma}[theorem]{Lemma}
\newtheorem{corollary}[theorem]{Corollary}
\newtheorem{assumption}[theorem]{Assumption}
\theoremstyle{definition}
\newtheorem{definition}[theorem]{Definition}
\newtheorem{remark}[theorem]{Remark}
\newcommand{\Q}{\mathbb Q}
\newcommand{\Z}{\mathbb Z}
\newcommand{\A}{\mathbb A}
\newcommand{\OO}{\mathcal O}
\newcommand{\Src}{\mathcal S}
\DeclareMathOperator{\Spec}{Spec}
\DeclareMathOperator{\Frac}{Frac}
\title{Variable bounds and exact transformations for a proposed\newline
Diophantine definition of the integers in the rationals}
\author{Research report prepared for JUVa\thanks{Prepared with ChatGPT and
parallel mathematical and computational reviews. No claim of independent
human review or priority is made. The upstream manuscript and its arithmetic
inputs are attributed separately.}}
\date{6 November 2026\quad---\quad working version for review}
\begin{document}
\maketitle
\begin{abstract}
We study polynomial presentations of the ordinary integers in the
rationals, starting from the repository \emph{diophantine-z-q} at a pinned
revision. Under its stated arithmetic inputs and an additional
simultaneous prime-pattern theorem of Sun, we construct a new polynomial
in 43 total variables, and a quartic in 1367 total variables. These are
different arithmetic presentations from the released 12,504-variable
quartic. A 45-variable alternative has substantially better degree bounds
than the 43-variable construction, and a fully expanded 1391-variable
quartic has coefficient height less than $2^{1020}$.

The small-variable polynomials are explicit indexed determinant and
finite-product formulas with proved degree and coefficient bounds.
Their enormous matrices and outer products are not expanded. The
1367-variable quartic has a complete coefficient table over one specified
fixed integer; the 1391-variable table consists of ordinary integer
coefficients. We also implement and independently check the
558-variable finite-algebra construction, including all exceptional
quadratic strata, rather than leaving it as an existence argument.

The lower bounds are independent of those arithmetic assumptions.
Any defining polynomial has degree at least four and at least two
existential witnesses. A quartic presented as a sum of squares of
quadratic polynomials needs at least five total variables. A nonnegative
polynomial in three total variables defining $\mathbb Z$ in $\mathbb Q$
needs degree at least eight.
Two witnesses cannot suffice when their joint degree is at most two,
irrespective of the degree in the input; the proof gives a simultaneous
conic-specialization obstruction. We also prove precise integer-to-rational
reductions and quantitative MRDP consequences. Finally, the released
polynomial accepts every input over $\Q(i)$, while the existence of
field-dependent definitions of all rings of integers remains unsettled
by this work. All hypotheses, exact checks, and remaining arithmetic
audit obligations are stated explicitly.
\end{abstract}
\tableofcontents

\section{Definitions, objectives, and the scope of the results}
\label{sec:scope}

\begin{definition}
Let $k$ be a field. A set $S\subseteq k$ is \emph{Diophantine over $k$}
if there are $r\geq0$ and $P\in k[T,Y_1,\ldots,Y_r]$ such that
\[
 t\in S\quad\Longleftrightarrow\quad
 \exists(y_1,\ldots,y_r)\in k^r:\ P(t,y_1,\ldots,y_r)=0.
\]
The variable $T$ is the input. We call $r$ the number of witnesses and
$r+1$ the total number of variables. Degrees in this report are total
degrees, including the input. For a nonzero polynomial over $\Q$, clear
numerical denominators and divide the coefficients by their greatest
common divisor. Its \emph{coefficient height} is the maximum absolute
value of the resulting primitive integer coefficients.
\end{definition}

The requested order of optimization is lexicographic: first the total
number of variables, then degree, then coefficient height. Under this
order, a smaller-variable polynomial of very high degree is preferable
to a quartic with more variables. A useful practical comparison should
nevertheless retain the alternatives with moderate coefficients and
complete coefficient tables. A circuit node used to \emph{write} a
polynomial is not a witness variable. Conversely, a newly quantified
coordinate constrained to equal a circuit value does count as a witness.

We use the upstream repository and the user's fork at revision
\begin{center}\small\ttfamily
28e87b1b62af0effd954605af7f70b2699dd5672.
\end{center}
The released package is \texttt{quartics/quartic-12504}; its manuscript
is \texttt{quartic\_12504.tex}. The repository's own documentation
distinguishes finite algebraic verification from verification of the
arithmetic theorem; its Lean material is incomplete
\cite{Repository2026}. We retain that distinction throughout.

Write $\Src$ for the rational existential system specified by the
manuscript's reciprocal comparison source, with precisions $64$ and
$16$, and its scalar order $48$. Write $S_{\Src}$ for its projected
rational solution set in the ordinary input $X$.

\begin{assumption}[The source arithmetic assertion]\label{ass:source}
The set $S_{\Src}$ is the ordinary subset $\Z\subset\Q$.
\end{assumption}

This is the assertion of the source manuscript's theorem on the
arithmetic interpretation of the revised source. It is not an additional
conjecture introduced by the optimizations. The source-preserving transformations
establish equality with $S_{\Src}$; calling their outputs definitions
of $\Z$ requires Assumption~\ref{ass:source}. Replaying a certificate
for polynomial substitution does not prove this assumption. In
particular, the present report does not claim a new independent proof
of the full arithmetic manuscript.

The new shared-point construction uses more information than the bare
equality in Assumption~\ref{ass:source}: it uses the manuscript's stated
cyclic-group, pole, formal-logarithm, height, and Lucas results, together
with Sun's simultaneous prime-pattern theorem
\cite[Theorem~4.1]{Sun2026Rational86}. Its soundness and completeness
are proved from these explicitly named inputs. This additional theorem
permits comparison points to be shared across different selected primes.
We distinguish this arithmetic modification from the unconditional
algebraic transformations of the released source. The lower-bound
theorems require neither arithmetic assumption.

The source release has 12,504 total variables, 11,939 quadratic
residuals, 135,721 monomials in their sum of squares, and maximum
coefficient
\[
 685783949153203534233601<2^{80}.
\]
These are exact algebraic statistics, independently replayed here.
The following results have different forms of explicitness.

\begin{center}
\small
\begin{tabular}{r l l p{43mm}}
\toprule
Total variables & Degree & Coefficients & Output supplied\\
\midrule
43 & $\le2^{9566680}$ & $\le2^{2^{86098389}}$ & Indexed polynomial, including a fixed elliptic exclusion product\\
45 & $\le2^{243}$ & $\le2^{2^{246}}$ & Indexed polynomial with materialized coefficient DAG\\
558 & $\le2^{3122}$ & $\le2^{2^{3127}}$ & Implemented indexed finite-algebra polynomial\\
1367 & $4$ & Fixed $\kappa=3^{2^{28}}$ & Complete table over $\mathbb Z[\kappa]$\\
1391 & $4$ & $<2^{1020}$ & Complete integer coefficient table\\
12,504 & $4$ & $<2^{80}$ & Unchanged released package\\
\bottomrule
\end{tabular}
\end{center}

All counts include the input. The first, second, fourth, and fifth rows
use the additional shared-point arithmetic argument. The third row is a
source-preserving algebraic construction. The displayed large degree and
height bounds are certified upper bounds, not exact expansion statistics.
Further source-preserving 853-variable circuits and 12,449--12,493-variable
quartics are retained as independently reviewable comparison constructions.
None of these upper bounds is claimed to be globally minimal. The
unrestricted gap is between three and 43 total variables under the stated
arithmetic inputs; for the supplied sum-of-squares quartic architecture
the gap is between five and 1367. Degree four is the least possible
degree if the conditional quartic construction is valid.

\section{From integer equations to rational equations}
\label{sec:reduction}

Suppose first that an integer-coefficient polynomial
$P(T,\mathbf Y)$ of degree $d$ defines $\Z$ in $\Q$, using $r$
witnesses. For $f\in\Z[X_1,\ldots,X_m]$, form
\begin{equation}\label{eq:basic-reduction}
 R_f(\mathbf X,\mathbf Y_1,\ldots,\mathbf Y_m)
 =f(\mathbf X)^2+\sum_{i=1}^m P(X_i,\mathbf Y_i)^2.
\end{equation}
\begin{proposition}\label{prop:integer-rational-reduction}
The equation $f=0$ has a solution in $\Z^m$ if and only if
$R_f=0$ has a solution over $\Q$. The construction is effective,
uses $m(r+1)$ rational variables, and has degree at most
$2\max\{\deg f,d\}$.
\end{proposition}
\begin{proof}
Given an integer zero of $f$, the defining property of $P$ supplies
a witness tuple for each coordinate, so every summand in
\eqref{eq:basic-reduction} vanishes. Conversely, a rational sum of
squares vanishes only when every summand vanishes. Thus any rational
zero of $R_f$ satisfies $f=0$ and $P(X_i,\mathbf Y_i)=0$ for every
$i$. Each $X_i$ is therefore an integer. The construction consists
only of copying, squaring, and adding finite polynomials; the bounds
follow by inspection.
\end{proof}

For example, an equation $f(a,b,c)=0$ is changed into a rational
equation with three separate integrality tests, one for each of
$a,b,c$. Its witnesses are allowed to have denominators. The tests
force the three original coordinates to be integers; they need not
force the added witnesses to be integers.

There is a better degree bound for the quartics actually supplied
here. They have the form $P=\sum_j q_j^2$, with $\deg q_j\leq2$.
In that case $P$ is nonnegative on rational inputs, so
\begin{equation}\label{eq:sos-reduction}
 f(\mathbf X)^2+\sum_{i=1}^m P(X_i,\mathbf Y_i)=0
\end{equation}
already gives the desired reduction. Its degree is at most
$\max\{2\deg f,4\}$, with the same variable count. Squaring each
copy of $P$ a second time is unnecessary.

Moreover every polynomial $f$ can be evaluated by a finite sequence
of additions and multiplications. Introduce one rational coordinate
for each intermediate value, and impose equations such as
$z-u-v=0$ or $z-uv=0$, together with the equation that the final value
is zero. These equations have degree at most two, and for fixed
$\mathbf X$ their intermediate coordinates are uniquely determined.
Combining their squares with the integrality quartics produces a
quartic rational-solvability problem for \emph{every} $f$. If $g$
intermediate coordinates are used, the variable count is
$m(r+1)+g$. These intermediate coordinates need no separate
integrality predicates: they already equal the prescribed polynomial
values. This trade preserves degree at the cost of variables, so it
is secondary under the requested priority.

\subsection{The precise contribution of MRDP}

The Matiyasevich--Robinson--Davis--Putnam theorem identifies the
computably enumerable subsets of $\Z^n$ with the integer Diophantine
sets, with effective and uniform polynomial representations
\cite{Matiyasevich1970,Jones1982}. A set is computably enumerable if
an algorithm can list its members, without necessarily deciding
nonmembership. The theorem over nonnegative integers gives the version
over all integers by the usual signed encodings and four-square
representations of nonnegative variables.

\begin{corollary}\label{cor:rational-mrdp}
If $\Z$ is Diophantine in $\Q$, then the computably enumerable
subsets of $\Q^n$ are exactly its Diophantine subsets. In particular,
rational polynomial solvability is computably enumerable complete
under many-one reductions, and hence undecidable.
\end{corollary}
\begin{proof}
Use the unique coprime numerator--positive denominator encoding of
each rational coordinate. If $S\subseteq\Q^n$ is computably
enumerable, its set of such integer encodings is computably enumerable.
MRDP supplies an integer polynomial for that set. Replace each integer
quantifier by a rational quantifier and an integrality test $P=0$;
add the equations $b_i x_i-a_i=0$. The conjunction is one polynomial
by a sum of squares, and has exactly the desired rational projection.
Conversely, enumerate all rational witness tuples and evaluate a given
polynomial exactly; this enumerates every Diophantine set.
Uniformity gives many-one hardness from any computably enumerable
set, and direct enumeration gives membership in that class.
\end{proof}

Applying this argument to a universal computably enumerable relation
also gives a universal rational Diophantine polynomial, with its index
as an input parameter. This does not make the small integer universal
polynomials into definitions over rational witnesses without modification.
Nor is an abstract Diophantine model of arithmetic automatically a
definition of the ordinary embedded subset $\Z\subset\Q$. A model,
an undecidability reduction, and a definition of that particular subset
are different assertions. All the conclusions in this section become
conditional applications to the source by
Assumption~\ref{ass:source}.

\section{The arithmetic structure retained from the source}
\label{sec:source-structure}

The polynomial is assembled from a considerably smaller list of
arithmetic requirements. It uses
\[
 K=\Q(\sqrt2),\quad d=3-\sqrt2,\quad D=d^2=11-6\sqrt2,
 \qquad E_D:y^2=x^3-Dx,
\]
and the point
\[
 P_*=\left(\frac{57}{8}-\frac72\sqrt2,
 -\frac{165}{8}+\frac{417}{32}\sqrt2\right).
\]
The source's generator assertion identifies $2E_D(K)$ with the
infinite cyclic group generated by $P_*$. Its points can therefore
encode integer indices. This assertion is a substantive arithmetic
input; the elementary chart substitutions proved below do not require
us to infer it from a computer search.

The input is converted to $\beta=X^2+2$. Four points encode indices
$n,s,t,z$, and three rational proxy coordinates satisfy a Lucas-type
equation. Formal-group comparisons at selected primes make the proxies
congruent to the indices. A common denominator and forty-one affine
contact expressions convert these congruences into height bounds.
A separate scalar batch bounds the height of $\beta$ in terms of
$|s|$. The anchor forces $|s|$ beyond the fixed cutoff. Finally Lucas
growth contradicts the height bounds unless the Lucas residual is zero
and $\beta=n\in\Z$. A rational number whose square is an integer
is an integer, so $X\in\Z$ follows.

In the other direction, the manuscript chooses sufficiently divisible
elliptic indices, uses a simultaneous ordinary-prime pattern, constructs
private comparison points, and obtains pole witnesses. The difficult
inputs include its quantitative five-contact estimate, its restricted
converse and pole-completeness argument, the prime-pattern theorem,
and the cyclic elliptic-group assertion. We preserve these inputs and
their precise witness scopes. Our changes operate on the resulting
existential presentation.

Two forms of verification should be separated. The released certificate
is a finite algebraic object: it checks lowering, elimination, invertible
row operations, variable maps, and expansion into a sum of squares.
The arithmetic manuscript must additionally prove both soundness and
completeness of its source. The former can be checked in seconds by exact
integer arithmetic. It cannot establish the latter by itself. The
proofs and evidence supplied below explicitly identify which task each
calculation performs.
